\section{Introduction}

Robot-learning evaluation increasingly depends on a chain of conversions: recorded episodes become
training datasets, reconstructed scenes become simulator assets, simulator rollouts become
benchmark results, and generated counterfactuals become data augmentation. The chain is only as
scientific as its hidden alignment assumptions. A camera transform with the wrong direction, an
action stream interpreted in the wrong frame, a command timestamp treated as an actuation timestamp,
or a train/test split that shares the same reconstructed world can change a conclusion without
changing a dataset's apparent file validity.

Existing formats solve important pieces of this chain. The LeRobot library vertically integrates
robot learning from motor interfaces through dataset collection, storage, streaming, policies, and
deployment~\citep{lerobot}. The current LeRobotDataset v3 layout stores high-frequency states,
actions, and timestamps in Parquet, visual streams in MP4 shards, and episode indexing in metadata
rather than filenames~\citep{lerobot_v3}. Rerun provides a physical-data layer and storage target
for multimodal recordings. NCore documents explicit pose graphs and sensor conventions. OpenUSD,
SimReady, glTF, and simulator-native formats describe composed worlds and assets. Gaussian splats
are also moving into standard asset containers through glTF's \texttt{KHR\_gaussian\_splatting}
extension~\citep{khronos_gs, gltf_gs}. Recent real-to-sim systems such as RoboSnap and DROID-Sim
show that Gaussian-splat environments can support trajectory replay, synthetic data generation, and
policy evaluation at dataset scale~\citep{robosnap}.

This landscape rules out a weak framing: WorldEpisode is not the first format that connects robot
data and Gaussian-splat worlds. That space is already occupied by asset standards and method papers.
The remaining opening is the contract connecting all of them: the persistent, auditable binding
between a robot-learning episode and the versioned world in which it occurred.

WorldEpisode targets that missing semantic layer. The contribution is a formal semantic contract,
not a wrapper around existing libraries. It is:
\begin{itemize}[leftmargin=*]
  \item \textbf{storage-neutral}: serializable through LeRobotDataset v3, Rerun, NCore, MCAP, or a reference package;
  \item \textbf{representation-neutral}: compatible with splats, meshes, NeRFs, point clouds, and future assets;
  \item \textbf{runtime-neutral}: transferable to Isaac Sim, MuJoCo, SAPIEN, Genesis, and other simulators;
  \item \textbf{loss-explicit}: conversions may be lossy, but never silently lossy.
\end{itemize}

Figure~\ref{fig:worldepisode-bridge} summarizes the intended relationship: WorldEpisode sits
between existing dataset/log containers and world/simulator assets as a semantic contract, not as a
replacement storage engine.

\begin{figure}[t]
\centering
\resizebox{\linewidth}{!}{%
\begin{tikzpicture}[font=\sffamily]
  \node[wecard=weBlue, text width=2.75cm, minimum height=1.0cm] (lerobot) at (0,1.8) {\textbf{LeRobot v3}\\state/action/video};
  \node[wecard=weCyan, text width=2.75cm, minimum height=1.0cm] (rerun) at (0,0.45) {\textbf{Rerun / MCAP}\\multirate logs};
  \node[wecard=weViolet, text width=2.75cm, minimum height=1.0cm] (ncore) at (0,-0.9) {\textbf{NCore}\\sensor calibration};

  \node[wecard=weInk, fill=wePanel, text width=4.05cm, minimum height=3.85cm] (core) at (5.0,0.45) {};
  \node[font=\sffamily\bfseries\large, text=weInk] at (core.north) [yshift=-0.35cm] {WorldEpisode};
  \node[font=\sffamily\small, text=weSlate] at (5.0,1.32) {storage-neutral semantic contract};
  \node[wechip=weBlue] at (3.85,0.65) {identity};
  \node[wechip=weCyan] at (6.05,0.65) {frames + clocks};
  \node[wechip=weAmber] at (3.85,-0.05) {action contract};
  \node[wechip=weGreen] at (6.05,-0.05) {roles + events};
  \node[wechip=weRose] at (4.85,-0.75) {provenance + lineage};

  \node[wecard=weGreen, text width=2.75cm, minimum height=1.0cm] (usd) at (10.0,1.8) {\textbf{OpenUSD}\\world composition};
  \node[wecard=weAmber, text width=2.75cm, minimum height=1.0cm] (gltf) at (10.0,0.45) {\textbf{glTF / SPZ}\\Gaussian splats};
  \node[wecard=weRose, text width=2.75cm, minimum height=1.0cm] (sim) at (10.0,-0.9) {\textbf{MuJoCo / Isaac}\\replay targets};

  \draw[wearrow=weBlue] (lerobot.east) -- ($(core.west)+(0,1.0)$);
  \draw[wearrow=weCyan] (rerun.east) -- (core.west);
  \draw[wearrow=weViolet] (ncore.east) -- ($(core.west)+(0,-1.0)$);
  \draw[wearrow=weGreen] ($(core.east)+(0,1.0)$) -- (usd.west);
  \draw[wearrow=weAmber] (core.east) -- (gltf.west);
  \draw[wearrow=weRose] ($(core.east)+(0,-1.0)$) -- (sim.west);

  \node[wechip=weInk, fill=wePanel, text width=9.6cm, align=center] at (5,-2.05)
    {Executable outputs: validator diagnostics \quad conversion-loss reports \quad lineage-safe splits \quad replay assumptions};
\end{tikzpicture}%
}
\caption{WorldEpisode is the semantic bridge between robot-learning containers, reconstructed
world assets, and simulator targets. It preserves the world--episode contract without replacing
the storage formats on either side.}
\label{fig:worldepisode-bridge}
\end{figure}

The paper's core claim is therefore not a new monolithic file format. It is an executable
interchange profile for world--episode semantics: identity, frames, clocks, action contracts,
representation roles, events, provenance, uncertainty, and lineage-safe evaluation splits. The
current artifact includes a JSON Schema, semantic validator, round-trip binding artifacts, an
active public LeRobotDataset v3 conversion experiment, hand-authored independent fixtures, and
controlled experiments for replay drift, leakage, and counterfactual augmentation.
